\chapter{Введение}

Актуальность работы обусловлена тем, что современные операционные системы используют модель управления памятю, сложившуюся с появлением UNIX. Она требует от программиста явного сохранения данных на диск, что увеличивает объём кода и приводит к потерям при сбоях~\cite{PLA}. Концепция персистентности устраняет эти недостатки, однако её практические реализации на современном оборудовании по-прежнему остаются предметом исследований~\cite{Revisited,Aurora}.

Работа выполнена в контексте проекта по портированию Phantom~OS на Genode. Phantom~OS реализует ортогональную персистентность внутри Phantom Virtual Machine~(PVM), периодически сохраняя её состояние на диск ~\cite{Phantom_docs}. Снапшоты хранятся в открытом виде, что создаёт угрозу конфиденциальности: в момент снимка в памяти могут быть секретные данные в открытом виде~\cite{SNIA_PM_Security,VM_Security}. \textbf{Цель работы}~--- обеспечить конфиденциальность снапшотов в состоянии покоя. Для этого адаптирован компонент, который реализует шифрование на блочном уровне, используя библиотеку Tresor, проведена верификация защиты и оценено влияние на производительность.
